\section{Hotword Conditioning}
\label{sec:hotwords}

\subsection{Setup}

The CommonVoice 17 hotword test splits ship with a per-utterance
ground-truth list of rare entities (mined offline) and a per-utterance
QA acceptance flag. For a target prompt size $k$, the inference client
pads each utterance's list to $k$ entries by mixing the ground-truth
hotwords with random distractors drawn from the same corpus's hotword
pool; no online hard negatives are added at evaluation time, in
contrast to training (Section~\ref{subsec:hotword-pool}). The $k=0$
row drops the hotword prompt line entirely (mirroring the training-time
fallback) and is identical to the corresponding row of
Tables~\ref{tab:asr_zh}--\ref{tab:asr_en}. We report CER for Mandarin
and WER for English.

Training-time hotword exposure for the v3 checkpoint follows
Section~\ref{subsec:hotword-pool}: each sample's prompt is built from
\textbf{real terms} (each ground-truth term independently dropped with
probability $0.05$), \textbf{up to 10 hard negatives} from the
char-overlap + bigram-Jaccard + pinyin-syllable retriever, and the
\textbf{remainder of $n_\text{distractor}\!\in\![1, 20]$ slots} from
the global pool. The hotword prompt line itself is kept with
probability $0.8$. Training and evaluation sweeps therefore cover the
same prompt-length regime, but evaluation is strictly easier because
it omits hard-negative pressure.

\subsection{Results}

\begin{table}[t]
  \caption{Effect of in-prompt hotwords on CommonVoice 17. CER (\%) for
    Mandarin and WER (\%) for English. Lower is better. $k$ is the
    number of hotwords supplied through the prompt.}
  \label{tab:hotwords}
  \centering
  \small
  \begin{tabular}{lrccccc}
    \toprule
    & & $k=0$ & $k=5$ & $k=10$ & $k=15$ & $k=20$ \\
    \midrule
    CV17 zh hotword test  & $n=10{,}647$ & 6.57 & \textbf{2.83} & 2.85 & 2.89 & 2.91 \\
    CV17 en hotword test  & $n=16{,}294$ & 9.65 & \textbf{4.89} & 4.99 & 5.03 & 5.09 \\
    \bottomrule
  \end{tabular}
\end{table}

\subsection{Analysis}

Two effects are visible in Table~\ref{tab:hotwords}:

\paragraph{Large absolute gains from a small list.}
Five hotwords already deliver a $3.74$ absolute CER reduction on
Mandarin (from $6.57$\% to $2.83$\%) and a $4.76$ absolute WER reduction
on English (from $9.65$\% to $4.89$\%). This is consistent with the
intuition that the rare-entity tail dominates the residual error of an
otherwise-well-trained ASR model, and that an LLM is well placed to
``snap'' to the suggested term when it is acoustically supported.

\paragraph{Diminishing -- and slightly inverted -- returns past five.}
Going from five to twenty hotwords has a small but consistent negative
effect (roughly $+0.08$ CER on Mandarin, $+0.20$ WER on English).
Because the eval client pads with \emph{random} distractors rather
than the hard negatives used during Stage~3 training, the residual
distractor pressure is mild but not absent; we attribute the slow
upturn to false-positive triggering on plausible-looking padding
terms. Section~\ref{sec:analysis} returns to this with a per-utterance
breakdown of insertion versus substitution errors.
% TODO: re-run the sweep with hard-negative padding to bound the
% adversarial worst case; current numbers reflect only random
% distractors.

\paragraph{Implication for deployment.}
For latency- and quality-sensitive deployments, supplying around five
likely hotwords per utterance (rather than dumping a full domain
glossary) appears to be a strong default for AmphionASR.
